\section{The holonomy angle $\theta_p$ as primitive object}
\label{sec:holonomy_primitive}

This section establishes the fundamental result that distinguishes
PT from classical gauge theories: under the structural hypotheses
of PT, the modular holonomy angle $\theta_p$ is the
\emph{unique} primitive invariant simultaneously satisfying four
natural conditions, and the classical pair $(A_\mu, F_{\mu\nu})$
is entirely reconstructible as a derivative of a more primitive
object --- the scalar persistence potential $S = -\ln \alpha$.
Position C of the foundational debate is no longer an
interpretation: it is a structural consequence.

\subsection{Proposition 3.1: $A_\mu$ and $F_{\mu\nu}$ are derived}
\label{ssec:prop_3_1_en}

\begin{proposition}[Persistence Potential Decomposition]
\label{prop:A_F_derived_en}
Let $S(\mu) = -\ln \alpha(\mu)$ be the persistence potential. Then
\begin{itemize}
  \item the 4-potential $A_\mu$ is identifiable with the gradient
  $\partial_\mu S$;
  \item the field tensor $F_{\mu\nu}$ is identifiable with the
  antisymmetric components of $\mathrm{Hess}(S)$;
  \item the metric $g_{\mu\nu}$ is identifiable with the symmetric
  components of $\mathrm{Hess}(S)$ (Lemma F, ch12 monograph).
\end{itemize}
\end{proposition}

\begin{proof}[Proof]
The construction is detailed in chapter~12 of the
monograph~\cite{PT_Monograph}; we recall here the three key steps.

\textbf{Step 1 (gradient identification).}
The Fisher metric associated with the stationary distribution
$\pi_m$ of the sieve is computed as the Hessian of the logarithm
of the partition function $Z_{\rm Ruelle}$, and the function
$\alpha(\mu)$ appears as the first eigenfunction of the transfer
operator. By construction, the gradient
$\partial_\mu S = -\partial_\mu \ln \alpha$ plays in the PT formalism
the role played by $A_\mu$ in classical electrodynamics: it enters
linearly in the interaction term of the effective Lagrangian
(\S~F3 of the monograph) and undergoes, under modular gauge
transformations $\theta_p \mapsto \theta_p + 2\pi k$, a phase
shift compatible with classical gauge invariance.

\textbf{Step 2 (antisymmetric Hessian identification).}
The antisymmetric components of the Hessian
$\mathrm{Hess}(S)_{\mu\nu} = \partial_\mu \partial_\nu S$ encode
the curvatures derived from the scalar potential, and reproduce
the six degrees of freedom of $F_{\mu\nu} = \partial_\mu A_\nu -
\partial_\nu A_\mu$ by direct identification. The identity
$dA = F$ becomes, under PT, the tautological identity
$d(dS) = 0$ projected onto the antisymmetric components of
$\mathrm{Hess}(S)$. This projection is non-trivial in the presence
of the Lorentzian signature $g_{00} < 0$, which separates the
spatial and temporal components of the Hessian.

\textbf{Step 3 (symmetric Hessian identification).}
The symmetric components of $\mathrm{Hess}(S)$ reproduce the
metric $g_{\mu\nu}$ of the emergent spacetime (Bianchi I on the
active primes $\{3, 5, 7\}$; see Lemma F of
ch.~12~\cite{PT_Monograph}). This identification is, in itself,
the content of the metric-reconstruction theorem: the
spacetime metric of PT \emph{is} the Hessian of the logarithm of
the partition function, with no new structure introduced.
\end{proof}

The immediate consequence is that $A_\mu, F_{\mu\nu}, g_{\mu\nu}$
are all three \emph{derivable from the same scalar potential}
$S(\mu)$, and that this potential is, in turn, computed from the
sole field $\{g_n\}$. None of these three objects is therefore
primary in the PT architecture: they are all three secondary
relative to $S$, and $S$ is itself secondary relative to
$\{g_n\}$. This structural observation underwrites the uniqueness
argument of the following subsection.

\subsection{Proposition 3.2: uniqueness of $\theta_p$ as primitive invariant}
\label{ssec:prop_3_2_en}

\begin{proposition}[Holonomic Primitiveness]
\label{prop:theta_p_unique_en}
The holonomy angle $\theta_p$ is the unique primitive invariant
simultaneously:
\begin{enumerate}[label=(\roman*),nosep]
  \item gauge-invariant,
  \item sensitive to Aharonov--Bohm,
  \item bounded,
  \item PT-definable from BA0 alone.
\end{enumerate}
\end{proposition}

\begin{proof}[Proof]
We show by exclusion that no other natural invariant
satisfies the four conditions simultaneously.

\textbf{Exclusion of $A_\mu$.}
The potential $A_\mu$ violates condition (i): it depends on gauge
by construction. No natural projection of $A_\mu$ is at once
gauge-invariant and bounded while remaining sensitive to
Aharonov--Bohm in the region where $F = 0$. Position A of the
foundational debate is therefore structurally incompatible with
the four conditions.

\textbf{Exclusion of $F_{\mu\nu}$.}
The field tensor $F_{\mu\nu}$ satisfies (i), (iii), (iv) but
violates condition (ii): it is, by construction, identically zero
in the interior of an ideal Aharonov--Bohm path, and so cannot
account for the observed phase shift. Position B of the debate
must add a supplementary non-local structure to recover sensitivity
to AB; that added structure is not PT-definable from BA0 alone.

\textbf{Exclusion of arbitrary functions of $\{g_n\}$.}
An arbitrary function $F[\{g_n\}]$ may satisfy (i), (ii), (iv) but
generally violates (iii): without a modular parallel-transport
structure, the function is not naturally bounded and does not
admit a simple identification with a measurable topological
invariant. Only functions that reduce to the modular angle
$\theta_p$ of a CRT channel $\mathbb{Z}/p\mathbb{Z}$ satisfy the
four conditions simultaneously.

\textbf{Existence and uniqueness.}
The angle $\theta_p$ defined by the
identity~\eqref{eq:sin2_theta_en} satisfies (i) (it is
gauge-invariant because defined as a closed integral),
(ii) (it is non-zero in every region where the modular connection
is non-trivial, including in the interior of an AB loop), (iii)
(it takes values in $[0, \pi]$ by construction), and (iv) (it is
defined by identity~\eqref{eq:sin2_theta_en}, which uses only
$\{g_n\}$ via $q(\mu)$). It is the unique primitive invariant
realising the four conditions.
\end{proof}

\subsection{The Wu--Yang--Healey position promoted to PT theorem}
\label{ssec:wu_yang_promoted_en}

Propositions~\ref{prop:A_F_derived_en}
and~\ref{prop:theta_p_unique_en} have a common reading: position
C of the foundational debate
(\S~\ref{ssec:position_C_en}) --- the primacy of holonomy --- is
no longer, in the PT framework, an interpretive choice among
others. It is a structural consequence: under the axioms of PT,
$\theta_p$ is the \emph{unique} primitive invariant available,
and $A_\mu, F_{\mu\nu}, g_{\mu\nu}$ are all three derived from
it via the scalar potential $S$.

This status has two immediate consequences. On one hand, PT
\emph{predicts} that any measurement of $A_\mu$ or $F_{\mu\nu}$
reduces, in its informational content, to a measurement of
$\theta_p$ on an active modular channel. No physical information
is lost by this reduction --- the AB phase is entirely encoded
in $\theta_3$ (cf.~\S~\ref{ssec:ab_consequence_en}). On the other
hand, PT \emph{excludes} the ontologies of position A (the reality
of the potential) and position B (the intrinsic non-locality of
the field): these ontologies postulate primitive objects that, in
the PT framework, cannot exist independently of the scalar
potential $S$ from which they are derived.

Position C becomes thus, under PT, more than an interpretation:
it is the only framework in which the four natural conditions
(i)--(iv) of Proposition~\ref{prop:theta_p_unique_en} can be
satisfied simultaneously. The foundational debate is, in this
precise measure, adjudicated.
